%---------------------------------------------------------------------
% Part two: Surah Aal-e-Imran. Six questions.
%
% Every verse quoted here was actually set in one of the four papers;
% none is chosen for illustration. References are surah:ayah and were
% checked individually.
%---------------------------------------------------------------------

\section*{حصہ دوم --- سورۃ آل عمران}
\markboth{سورۃ آل عمران}{}

\begin{nukta}[پہلے سورت کا نقشہ ذہن میں بٹھا لیجیے]
\textbf{مدنی} سورت، قرآن کی \textbf{تیسری} سورت، \textbf{دو سو} آیات، بیس
رکوع، تیسرے اور چوتھے پارے میں۔ نام حضرت مریمؑ کے خاندان ---
\emph{آلِ عمران} --- کے ذکر پر رکھا گیا (آیت \textenglish{33})۔

\smallskip
\textbf{مرکزی مضامین:} توحید و صفاتِ باری تعالیٰ \ \textbullet\ \ اہلِ کتاب سے
مکالمہ اور نجران کے وفد کا مباہلہ \ \textbullet\ \ حضرت مریمؑ و حضرت عیسیٰؑ کی
حقیقت \ \textbullet\ \ غزوہ احد کا تفصیلی تجزیہ (آیات
\textenglish{121--179}) \ \textbullet\ \ وحدتِ امت اور تفرقے کی ممانعت \
\textbullet\ \ امتِ مسلمہ کا منصب اور آخرت کی حقیقت۔

\smallskip
یہی چھ عنوان پورے حصے کی کنجی ہیں --- ہر موضوعاتی سوال انہی میں سے کسی ایک پر
ہوتا ہے۔
\end{nukta}

%---------------------------------------------------------------------
\begin{sawal}{سوال \textenglish{1} \ \textnormal{\small(چاروں پرچوں میں --- بہار \textenglish{2012} س\,\textenglish{1}، بہار \textenglish{2013} س\,\textenglish{2}، خزاں \textenglish{2021} س\,\textenglish{2}، بہار \textenglish{2022} س\,\textenglish{2})}}
سورۃ آل عمران کی درج ذیل آیات کا ترجمہ اور تشریح لکھیے۔
\end{sawal}

\begin{hal}
یہ سوال \textbf{ہر پرچے میں} آیا ہے۔ ذیل کی نو آیات وہی ہیں جو ان چاروں پرچوں
میں دی گئیں --- یعنی یہ محض متوقع نہیں، \emph{آزمودہ} ہیں۔

\ayat{إِنَّ الدِّينَ عِندَ اللهِ الْإِسْلَامُ}{آل عمران، \textenglish{19} ---
[خزاں \textenglish{2021}]}
\tarjuma{بے شک اللہ کے نزدیک دین صرف اسلام ہے۔}
\textbf{تشریح:} اسلام کوئی نیا مذہب نہیں، تمام انبیاء کا مشترک دین ہے؛ اللہ کے
ہاں قبولیت کا معیار صرف ایک۔ آگے آیت \textenglish{85} اسے دو ٹوک کر دیتی ہے:
\ar{وَمَن يَبْتَغِ غَيْرَ الْإِسْلَامِ دِينًا فَلَن يُقْبَلَ مِنْهُ}۔

\medskip
\ayat{لَا يَتَّخِذِ الْمُؤْمِنُونَ الْكَافِرِينَ أَوْلِيَاءَ مِن دُونِ
الْمُؤْمِنِينَ}{آل عمران، \textenglish{28} --- [خزاں \textenglish{2021}]}
\tarjuma{مومن، مومنوں کو چھوڑ کر کافروں کو اپنا رفیق و کارساز نہ بنائیں۔}
\textbf{تشریح:} ''ولی'' یعنی \emph{دلی دوستی اور وفاداری}، عام حسنِ سلوک نہیں
--- انصاف اور تجارت کی ممانعت نہیں، دین پر وفاداری کی بات ہے۔

\medskip
\ayat{قُلْ إِن كُنتُمْ تُحِبُّونَ اللهَ فَاتَّبِعُونِي يُحْبِبْكُمُ اللهُ
وَيَغْفِرْ لَكُمْ ذُنُوبَكُمْ وَاللهُ غَفُورٌ رَحِيمٌ}{آل عمران،
\textenglish{31} --- [بہار \textenglish{2013}]}
\tarjuma{کہہ دیجیے: اگر تم اللہ سے محبت رکھتے ہو تو میری پیروی کرو، اللہ تم سے
محبت کرے گا اور تمہارے گناہ بخش دے گا۔ اور اللہ بہت بخشنے والا، نہایت مہربان
ہے۔}
\textbf{تشریح:} یہ ''آیتِ محبت'' کہلاتی ہے۔ اس میں محبت کا \textbf{ثبوت} مقرر
کر دیا گیا: دعویٰ نہیں، \emph{اتباع}۔ اور صلہ بھی دو ہیں --- اللہ کی محبت اور
مغفرت۔ اسی آیت سے اتباع و اطاعت کا فرق بھی نکلتا ہے (دیکھیے حصہ اول)۔

\medskip
\ayat{وَاعْتَصِمُوا بِحَبْلِ اللهِ جَمِيعًا وَلَا تَفَرَّقُوا وَاذْكُرُوا
نِعْمَتَ اللهِ عَلَيْكُمْ إِذْ كُنتُمْ أَعْدَاءً فَأَلَّفَ بَيْنَ قُلُوبِكُمْ
فَأَصْبَحْتُم بِنِعْمَتِهِ إِخْوَانًا}{آل عمران، \textenglish{103} --- [بہار
\textenglish{2022}]}
\tarjuma{اور سب مل کر اللہ کی رسی کو مضبوطی سے تھام لو اور تفرقے میں نہ پڑو۔
اور اپنے اوپر اللہ کے اس احسان کو یاد کرو جب تم ایک دوسرے کے دشمن تھے، پھر اس
نے تمہارے دلوں کو جوڑ دیا اور تم اس کے فضل سے بھائی بھائی بن گئے۔}
\textbf{تشریح:} ''حبل اللہ'' سے مراد قرآن اور دینِ اسلام ہے۔ لفظ
\ar{جَمِيعًا} پر غور کیجیے --- حکم اجتماعی ہے، انفرادی نہیں۔ اور دلیل تاریخی
دی گئی: اوس و خزرج صدیوں سے برسرِ پیکار تھے، اسلام نے انہیں بھائی بنا دیا۔
آیت آگے چلتی ہے کہ تم آگ کے گڑھے کے کنارے تھے، اللہ نے بچا لیا۔

\medskip
\ayat{كُنتُمْ خَيْرَ أُمَّةٍ أُخْرِجَتْ لِلنَّاسِ تَأْمُرُونَ
بِالْمَعْرُوفِ وَتَنْهَوْنَ عَنِ الْمُنكَرِ وَتُؤْمِنُونَ
بِاللهِ}{آل عمران، \textenglish{110} --- [بہار \textenglish{2013}]}
\tarjuma{تم بہترین امت ہو جو لوگوں کے لیے نکالی گئی ہے؛ تم نیکی کا حکم دیتے
ہو، برائی سے روکتے ہو اور اللہ پر ایمان رکھتے ہو۔}
\textbf{تشریح:} یہ امتِ مسلمہ کا \emph{منصب نامہ} ہے۔ خیریت کوئی نسلی اعزاز
نہیں --- آیت نے اس کے ساتھ ہی تین شرطیں لگا دیں: امر بالمعروف، نہی عن المنکر،
اور ایمان باللہ۔ جو امت یہ تینوں چھوڑ دے، وہ اس آیت کی مصداق نہیں رہتی۔ لفظ
\ar{لِلنَّاسِ} بتاتا ہے کہ یہ منصب پوری انسانیت کے لیے ہے، صرف اپنوں کے لیے
نہیں۔

\medskip
\textbf{بہار \textenglish{2012} کے پرچے کی چار مزید مختصر آیات:}

\ayat{وَلِلهِ مِيرَاثُ السَّمَاوَاتِ وَالْأَرْضِ}{آل عمران، \textenglish{180}}
\tarjuma{اور آسمانوں اور زمین کی میراث اللہ ہی کی ہے۔}
\textbf{تشریح:} بخل کے بیان میں --- روکا ہوا مال بھی بالآخر لوٹتا اسی کی طرف ہے۔

\ayat{وَمَا الْحَيَاةُ الدُّنْيَا إِلَّا مَتَاعُ الْغُرُورِ}{آل عمران،
\textenglish{185}}
\tarjuma{اور دنیا کی زندگی دھوکے کے سامان کے سوا کچھ نہیں۔}
\textbf{تشریح:} دنیا کی مذمت نہیں، اس کے دھوکے کی نشان دہی --- عارضی کو دائمی
سمجھ لینا۔

\ayat{وَإِن تَصْبِرُوا وَتَتَّقُوا فَإِنَّ ذَلِكَ مِنْ عَزْمِ
الْأُمُورِ}{آل عمران، \textenglish{186}}
\tarjuma{اور اگر تم صبر کرو اور تقویٰ اختیار کرو تو یہ بڑے حوصلے کے کاموں میں
سے ہے۔}
\textbf{تشریح:} مال و جان کی آزمائش اور تکلیف دہ باتوں کا جواب --- صبر اور
تقویٰ۔

\ayat{وَاللهُ عِندَهُ حُسْنُ الثَّوَابِ}{آل عمران، \textenglish{195}}
\tarjuma{اور اللہ ہی کے پاس بہترین اجر ہے۔}
\textbf{تشریح:} ہجرت اور جہاد کرنے والوں کے اجر کا بیان، اور صاف حکم کہ عمل کے
اجر میں \ar{مِّن ذَكَرٍ أَوْ أُنثَىٰ} کوئی فرق نہیں۔
\end{hal}

\begin{nukta}[آیت لکھنے کا طریقہ]
بہار \textenglish{2012} کے پرچے میں سوال کے الفاظ تھے: ''آیات \textbf{خوش خط
لکھ کر} ترجمہ بھی تحریر کریں۔'' یعنی خطاطی خود نمبروں کا حصہ ہے۔ آیت
\textbf{الگ سطر} میں، \textbf{اعراب کے ساتھ} اور \textbf{بیچ میں} لکھیے، پھر
نیچے ترجمہ، پھر تشریح۔ حوالہ (سورت اور آیت نمبر) لکھنا نہ بھولیے۔
\end{nukta}

%---------------------------------------------------------------------
\begin{sawal}{سوال \textenglish{2} \ \textnormal{\small(بہار \textenglish{2022}، سوال \textenglish{1})}}
سورۃ آل عمران کی روشنی میں توحیدِ باری تعالیٰ پر جامع نوٹ تحریر کریں۔
\end{sawal}

\begin{hal}
\textbf{تمہید۔} سورت کا آغاز ہی توحید سے ہوتا ہے: \ar{اللهُ لَا إِلٰهَ إِلَّا
هُوَ الْحَيُّ الْقَيُّومُ} (آیت \textenglish{2})۔ چونکہ یہ سورت اہلِ کتاب کے
عقائد کے جواب میں نازل ہوئی، اس کی توحید \emph{استدلالی} ہے۔

\medskip
\textbf{(۱) فی الذات} --- شریک، بیٹا یا ہمسر نہیں: \ar{شَهِدَ اللهُ أَنَّهُ
لَا إِلٰهَ إِلَّا هُوَ} (آیت \textenglish{18})۔ \quad
\textbf{(۲) فی الصفات} --- حیّ، قیّوم، عزیز، حکیم: \ar{لَا يَخْفَىٰ عَلَيْهِ
شَيْءٌ} (آیت \textenglish{5})۔

\smallskip
\textbf{(۳) فی الافعال} --- عزت و ذلت، بادشاہی سب اسی کے ہاتھ: \ar{تُؤْتِي
الْمُلْكَ مَن تَشَاءُ وَتَنزِعُ الْمُلْكَ مِمَّن تَشَاءُ} (آیت
\textenglish{26})۔ \quad
\textbf{(۴) فی العبادت} --- عبادت و اطاعت صرف اسی کا حق: \ar{وَلَا يَتَّخِذَ
بَعْضُنَا بَعْضًا أَرْبَابًا مِّن دُونِ اللهِ} (آیت \textenglish{64})۔

\medskip
\textbf{شرک کی تردید۔} عیسیٰؑ کی الوہیت کا دعویٰ اس دلیل سے رد: پیدائش حضرت
آدمؑ سے زیادہ غیر معمولی نہیں --- \ar{إِنَّ مَثَلَ عِيسَىٰ عِندَ اللهِ
كَمَثَلِ آدَمَ خَلَقَهُ مِن تُرَابٍ} (آیت \textenglish{59})۔ جو خود پیدا کیا
گیا ہو، وہ خالق نہیں ہو سکتا۔

\medskip
\textbf{خلاصہ۔} آل عمران کی توحید فلسفیانہ نہیں، \emph{عملی} ہے --- انسان کسی
اور کے آگے نہ جھکے، مال و اقتدار کو حقیقی طاقت نہ سمجھے، اطاعت صرف ایک کی کرے۔
\end{hal}

%---------------------------------------------------------------------
\begin{sawal}{سوال \textenglish{3} \ \textnormal{\small(بہار \textenglish{2012}، سوال \textenglish{2})}}
سورۃ آل عمران کی روشنی میں ''اسلام دینِ کائنات اور دینِ انسانیت'' کے موضوع پر
جامع نوٹ لکھیے۔
\end{sawal}

\begin{hal}
جواب کو دو حصوں میں تقسیم کیجیے --- یہی سب سے صاف طریقہ ہے۔

\medskip
\textbf{(الف) اسلام دینِ کائنات ہے۔} ہر ذرہ خالق کے بنائے ضابطے کا پابند ہے:

\ayat{وَلَهُ أَسْلَمَ مَن فِي السَّمَاوَاتِ وَالْأَرْضِ طَوْعًا وَكَرْهًا
وَإِلَيْهِ يُرْجَعُونَ}{آل عمران، \textenglish{83}}
\tarjuma{اور آسمانوں اور زمین میں جو کوئی ہے سب چار و ناچار اسی کے فرمانبردار
ہیں، اور اسی کی طرف لوٹائے جائیں گے۔}

\smallskip
سورج، چاند، ہوا سب \emph{کرہاً} (بے اختیار) تابع ہیں؛ انسان کو اختیار دیا گیا
کہ \emph{طوعاً} (اپنی مرضی سے) تابع ہو --- یہی اس کا امتحان ہے۔

\medskip
\textbf{(ب) اسلام دینِ انسانیت ہے۔} کسی قوم یا نسل کا دین نہیں:

\begin{itemize}\setlength{\itemsep}{1pt}
  \item \ar{إِنَّ الدِّينَ عِندَ اللهِ الْإِسْلَامُ} (\textenglish{19}) ---
        معیار سب کے لیے ایک۔
  \item \ar{كُنتُمْ خَيْرَ أُمَّةٍ أُخْرِجَتْ لِلنَّاسِ} (\textenglish{110})
        --- امت ''لوگوں کے لیے''، اپنے لیے نہیں۔
  \item تمام انبیاء کا دین ایک ہی بتایا گیا (\textenglish{84}) --- ابراہیم،
        موسیٰ، عیسیٰ سب کے درمیان فرق نہیں کیا جاتا۔
\end{itemize}

\medskip
\textbf{انسانی تقاضے۔} انصاف (\textenglish{18})، وحدت اور تفرقے کی ممانعت
(\textenglish{103})، مشورہ و نرمی (\textenglish{159})، مرد و عورت کی برابری
(\textenglish{195})۔

\medskip
\textbf{خلاصہ۔} کائنات پہلے ہی مسلم ہے؛ انسان کو دعوت ہے کہ اپنے اختیار سے اسی
صف میں شامل ہو جائے۔
\end{hal}

%---------------------------------------------------------------------
\begin{sawal}{سوال \textenglish{4} \ \textnormal{\small(بہار \textenglish{2013}، سوال \textenglish{7} اور خزاں \textenglish{2021}، سوال \textenglish{4})}}
سورۃ آل عمران کی روشنی میں تفرقے کی مذمت اور اہلِ کتاب کی خرابیوں پر نوٹ لکھیے۔
\end{sawal}

\begin{hal}
\textbf{(الف) تفرقہ --- ممانعت، سبب اور علاج۔} حکم صریح ہے: \ar{وَاعْتَصِمُوا
بِحَبْلِ اللهِ جَمِيعًا وَلَا تَفَرَّقُوا} (آیت \textenglish{103})، اور اس کے
فوراً بعد اہلِ کتاب کی مثال دی گئی:

\ayat{وَلَا تَكُونُوا كَالَّذِينَ تَفَرَّقُوا وَاخْتَلَفُوا مِن بَعْدِ مَا
جَاءَهُمُ الْبَيِّنَاتُ}{آل عمران، \textenglish{105}}
\tarjuma{اور ان لوگوں کی طرح نہ ہو جاؤ جو واضح دلائل آ جانے کے بعد تفرقے میں
پڑ گئے اور اختلاف کرنے لگے۔}

\smallskip
غور کیجیے: مذمت اس تفرقے کی ہے جو \emph{علم آ جانے کے بعد} ہو --- یعنی جس کی
جڑ لاعلمی نہیں، بلکہ ضد، حسد اور دنیا کی طلب ہے۔ \textbf{علاج} بھی اسی سورت
میں ہے: کتاب پر اجتماع (\textenglish{103})، دعوت و احتساب کا منظم ادارہ
(\textenglish{104})، اور مشورہ، نرمی و عفو (\textenglish{159})۔

\medskip
\textbf{(ب) اہلِ کتاب کی خرابیاں۔} سورت نے یہ فہرست خود گنوائی ہے:

\begin{itemize}\setlength{\itemsep}{1pt}
  \item \textbf{کتمانِ حق اور تلبیس۔} \ar{لِمَ تَلْبِسُونَ الْحَقَّ
        بِالْبَاطِلِ وَتَكْتُمُونَ الْحَقَّ} (\textenglish{71}) --- جانتے
        بوجھتے حق چھپانا۔
  \item \textbf{تحریف۔} \ar{يَلْوُونَ أَلْسِنَتَهُم بِالْكِتَابِ}
        (\textenglish{78}) --- پڑھتے وقت زبان مروڑنا تاکہ سننے والا اسے کتاب
        کا حصہ سمجھے۔
  \item \textbf{خیانتِ امانت اور منظم سازش۔} (\textenglish{72, 75}) ---
        عذر: \ar{لَيْسَ عَلَيْنَا فِي الْأُمِّيِّينَ سَبِيلٌ}؛ اور صبح ایمان
        لا کر شام کو پھر جانا، تاکہ مسلمان شک میں پڑیں۔
  \item \textbf{عہد شکنی۔} (\textenglish{187}) --- میثاق کتاب واضح کرنے کا
        تھا؛ \ar{فَنَبَذُوهُ وَرَاءَ ظُهُورِهِمْ} --- پسِ پشت ڈال دیا۔
  \item \textbf{قتلِ انبیاء، اور بغیر کیے تعریف کی خواہش۔}
        (\textenglish{112, 181, 188})۔
\end{itemize}

\begin{nukta}[یہ نقطہ لکھنا جواب کو مکمل کرتا ہے]
سورت نے سب کو ایک لاٹھی سے نہیں ہانکا: \ar{لَيْسُوا سَوَاءً مِّنْ أَهْلِ
الْكِتَابِ أُمَّةٌ قَائِمَةٌ يَتْلُونَ آيَاتِ اللهِ آنَاءَ اللَّيْلِ}
(\textenglish{113}) --- ''یہ سب ایک جیسے نہیں۔'' جو طالب علم یہ آیت لکھ دیتا
ہے، وہ دکھا دیتا ہے کہ اس نے پوری سورت پڑھی ہے، صرف مذمت کی آیات نہیں۔
\end{nukta}

\medskip
\textbf{(ج) نجران کا وفد اور مباہلہ۔} سورت کا نام حضرت مریم\,علیہا السلام کے
خاندان پر ہے۔ عیسیٰ\,علیہ السلام کی الوہیت کا دعویٰ ایک جملے میں رد ہوا: بغیر
باپ کے پیدا ہونا خدائی کی دلیل نہیں، ورنہ آدم\,علیہ السلام کا دعویٰ زیادہ
مضبوط ہوتا کہ ان کا نہ باپ تھا نہ ماں --- \ar{إِنَّ مَثَلَ عِيسَىٰ عِندَ اللهِ
كَمَثَلِ آدَمَ خَلَقَهُ مِن تُرَابٍ} (\textenglish{59})۔ جب دلیل سے بات طے نہ
ہوئی تو \textbf{مباہلے} کی دعوت دی گئی (\textenglish{61}): دونوں فریق اہل و
عیال سمیت جمع ہوں اور جھوٹے پر اللہ کی لعنت مانگیں۔ نجران کے عیسائی سامنے آنے
سے کترا گئے اور جزیے پر راضی ہو گئے --- یہی دعوت کی سچائی کا عملی ثبوت بنا۔
\end{hal}

%---------------------------------------------------------------------
\begin{sawal}{سوال \textenglish{5} \ \textnormal{\small(بہار \textenglish{2022}، سوال \textenglish{3} اور سوال \textenglish{4} --- ایک ہی پرچے میں دو سوال)}}
سورۃ آل عمران کی روشنی میں غزوہ احد کے اسباب، واقعات اور قرآنی تجزیے پر تفصیلی
نوٹ لکھیں۔
\end{sawal}

\begin{hal}
بہار \textenglish{2022} میں \emph{دو} سوال اسی ایک غزوے پر تھے --- ایک اسباب
پر، ایک واقعات پر --- کیونکہ سورۃ آل عمران کی آیات \textenglish{121} سے
\textenglish{179}، یعنی تقریباً ایک تہائی سورت، اسی غزوے کا تجزیہ ہے۔

\medskip
\textbf{بنیادی معلومات۔} شوال \textenglish{3}ھ، جبلِ احد۔ قریش تین ہزار
(قیادت ابو سفیان)، مسلمان ایک ہزار --- تین سو عبد اللہ بن ابی کے ساتھ لوٹ
گئے، میدان میں صرف \textbf{سات سو} رہ گئے۔

\medskip
\textbf{(الف) اسباب:}
\begin{itemize}\setlength{\itemsep}{1pt}
  \item \textbf{بدر کا انتقام} --- قریش کے ستر سردار بدر میں مارے گئے تھے۔
  \item \textbf{تجارتی راستے کی حفاظت} --- شام کا راستہ مدینے کے قریب سے گزرتا
        تھا اور قریش کی معیشت اسی پر تھی۔
  \item \textbf{ساکھ کی بحالی} --- بدر کی شکست نے پورے عرب میں قریش کی سرداری
        مشکوک کر دی تھی۔
  \item \textbf{یہود اور منافقین کی خفیہ اعانت} --- جو مدینے کے اندر سے قریش
        کی حوصلہ افزائی کر رہے تھے۔
\end{itemize}

\medskip
\textbf{(ب) واقعات --- ترتیب کے ساتھ:}
\begin{enumerate}\setlength{\itemsep}{1pt}
  \item \textbf{مشورہ۔} نبی کریم \saw کی اپنی رائے مدینے میں رہ کر دفاع کی تھی،
        مگر پُرجوش نوجوانوں کی اکثریت باہر نکلنے کے حق میں تھی۔ آپ \saw نے
        \emph{اکثریت کی رائے} قبول فرمائی --- شوریٰ کا عملی نمونہ۔ راستے سے
        عبد اللہ بن ابی تین سو ساتھیوں سمیت واپس چلا گیا۔
  \item \textbf{درّے پر تیر انداز۔} پچاس تیر انداز حضرت عبد اللہ بن جبیرؓ کی
        قیادت میں مقرر ہوئے، حکم دو ٹوک تھا: \emph{فتح ہو یا شکست، جگہ نہ
        چھوڑنا۔}
  \item \textbf{ابتدائی فتح، پھر فیصلہ کن غلطی۔} قریش کے قدم اکھڑ گئے اور
        مسلمان مالِ غنیمت تک پہنچ گئے؛ اکثر تیر انداز یہ سمجھ کر کہ جنگ ختم ہو
        گئی، درّہ چھوڑ کر غنیمت سمیٹنے لگے۔
  \item \textbf{پلٹا۔} خالد بن ولید --- جو اس وقت مسلمان نہیں ہوئے تھے ---
        نے خالی درّے سے پیچھے آ کر حملہ کیا اور مسلمان دو پاٹوں میں آ گئے۔
  \item \textbf{نقصان اور افواہ۔} ستر صحابہ شہید ہوئے، جن میں حضرت حمزہؓ بھی
        تھے؛ نبی کریم \saw زخمی ہوئے۔ آپ \saw کی شہادت کی جھوٹی خبر پھیلی، جس
        پر آیت اتری: \ar{وَمَا مُحَمَّدٌ إِلَّا رَسُولٌ قَدْ خَلَتْ مِن
        قَبْلِهِ الرُّسُلُ} (\textenglish{144})۔
\end{enumerate}

\medskip
\textbf{(ج) قرآنی تجزیہ --- یہی حصہ سب سے زیادہ نمبر دلاتا ہے۔} قرآن نے شکست کے
اسباب خود گنوائے:

\ayat{وَلَقَدْ صَدَقَكُمُ اللهُ وَعْدَهُ إِذْ تَحُسُّونَهُم بِإِذْنِهِ حَتَّىٰ
إِذَا فَشِلْتُمْ وَتَنَازَعْتُمْ فِي الْأَمْرِ وَعَصَيْتُم مِّن بَعْدِ مَا
أَرَاكُم مَّا تُحِبُّونَ}{آل عمران، \textenglish{152}}
\tarjuma{اور اللہ نے تم سے اپنا وعدہ سچا کر دکھایا جب تم اس کے حکم سے انہیں
کاٹ رہے تھے، یہاں تک کہ جب تم نے کمزوری دکھائی، حکم میں جھگڑا کیا اور نافرمانی
کی --- اس کے بعد کہ اللہ نے تمہیں وہ دکھا دیا جو تم چاہتے تھے۔}

\smallskip
تین اسباب صاف لکھے ہیں: \textbf{فشل} (کمزوری)، \textbf{تنازع} (باہمی جھگڑا) اور
\textbf{عصیان} (حکم کی خلاف ورزی)۔ اور ذمہ داری بھی طے کر دی گئی: \ar{قُلْ هُوَ
مِنْ عِندِ أَنفُسِكُمْ} --- ''یہ تمہارے اپنے سبب سے ہے'' (\textenglish{165})۔

\smallskip
\textbf{حوصلہ بھی دیا گیا:} \ar{وَلَا تَهِنُوا وَلَا تَحْزَنُوا وَأَنتُمُ
الْأَعْلَوْنَ إِن كُنتُم مُّؤْمِنِينَ} (\textenglish{139})، اور شہداء کے بارے
میں: \ar{بَلْ أَحْيَاءٌ عِندَ رَبِّهِمْ يُرْزَقُونَ} (\textenglish{169})۔ اور
غلطی کرنے والوں کے لیے معافی و مشورہ: \ar{فَبِمَا رَحْمَةٍ مِّنَ اللهِ لِنتَ
لَهُمْ} (\textenglish{159})۔

\medskip
\textbf{اسباق۔} حکم کی خلاف ورزی فتح کو شکست میں بدل دیتی ہے \ \textbullet\ \
آزمائش سے مخلص اور منافق کی پہچان ہوتی ہے \ \textbullet\ \ قیادت کا حسنِ اخلاق
شکست کے بعد جماعت کو بکھرنے سے بچاتا ہے۔
\end{hal}
